\sekcja{Tłumaczenia\footnotemark}{}
%---------------------
\small
\vspace{-5mm}
\red{\textbf{Ant.1}} A Jakób * zrodził Józefa, męża Marji, z której się narodził Jezus, którego zowią Chrystusem.\\
\red{\textbf{Psalm 109}}\\
1. Rzekł Pan Panu memu: * Siądź po prawicy mojej:\\
2. Aż położę nieprzyjaciół twoich * podnóżkiem nóg twoich.\\
3. Berło mocy twojej wypuści Pan z Syjonu: * panuj wpośród nieprzyjaciół twoich.\\
4. Przy tobie panowanie w dzień mocy twej w blaskach świętości: * z żywota przed jutrzenką zrodziłem cię.\\
5. Przysiągł Pan i nic będzie żałował: * Ty jesteś kapłanem na wieki, według porządku Melchizedecha.\\
6. Pan po prawicy twojej, * poraził w dzień gniewu swego królów.\\
7. Będzie sądził narody, dopełni zagłady, * potłucze głowy wielu na ziemi.\\
8. Ze strumienia na drodze pić będzie: * dlatego wywyższy głowę.\\
%---------------------
\red{\textbf{Ant.2}} Posłan jest Anioł * Gabrjel od Boga do panny poślubionej mężowi, któremu było Józef, z domu Dawidowego: a imię panny Marya.\\
\red{\textbf{Psalm 110}}\\
1. Będę Cię wysławiał, Panie wszystkim sercem moim, * w radzie sprawiedliwych i w zgromadzeniu.\\
2. Wielkie są dzieła Pańskie, * doświadczone według wszystkiego upodobania Jego.\\
3. Chwalebne i wspaniałe jest dzieło Jego, * a sprawiedliwość Jego trwa na wieki wieków.\\
4. Uczynił pamiątkę dziwów swoich miłosierny i litościwy Pan: * dał pokarm tym, którzy się Go boją.\\
5. Pamiętać będzie wiecznie na przymierze swoje: * moc uczynków swoich oznajmi ludowi swemu:\\
6. Aby im dać dziedzictwo narodów: * dzieła rąk Jego są prawdą i sądem.\\
7. Pewne są wszystkie przykazania Jego, utwierdzone na wieki wieków, * uczynione w prawdzie i w prawości.\\
8. Zesłał odkupienie ludowi swemu, * ustanowił na wieki przymierze swoje.\\
9. Święte i straszne jest imię Jego: * początkiem mądrości jest bojaźń Pańska.\\
10. Zrozumienie dobre mają, wszyscy, którzy tak czynią: * chwała Jego trwa na wieki wieków.\\
%---------------------
\red{\textbf{Ant.3}} Gdy była poślubiona * Matka Jezusa, Marja Józefowi, pierwej niźli się zeszli, znaleziona jest w żywocie mająca z Ducha Świętego.\\
\red{\textbf{Psalm 111}}\\
1. Błogosławiony mąż, który się boi Pana: * w przykazaniach Jego będzie rozmiłowany wielce.\\
2. Możne na ziemi będzie potomstwo jego: * ród prawych będzie błogosławiony.\\
3. Sława i bogactwo w domu jego, * a sprawiedliwość jego trwa na wieki wieków.\\
4. Wzeszła w ciemnościach światłość prawym, * miłościwy, miłosierny i sprawiedliwy.\\
5. Miły jest człowiek, który się lituje i pożycza, będzie miarkował mowy swe z rozsądkiem, * bo na wieki nie będzie poruszony.\\
6. W wiecznej pamięci będzie sprawiedliwy: * nie będzie się bał pogłosek złych.\\
7. Gotowe serce jego nadzieję mieć w Panu, umocnione jest serce jego: * nie poruszy się, aż wzgardzi nieprzyjaciółmi swymi.\\
8. Rozproszył, dał ubogim, sprawiedliwość jego trwa na wieki wieków, * róg jego podniesie się w chwale.\\
9. Grzesznik ujrzy i będzie się gniewał, będzie zgrzytał zębami swymi i będzie schnął, * żądza niezbożnych zaginie.\\
%---------------------
\red{\textbf{Ant.4}}  Józef, mąż jej, * będąc sprawiedliwym i nie chcąc jej osławiać, chciał ją potajemnie opuścić.\\
\red{\textbf{Psalm 112}}\\
1. Chwalcie, słudzy, Pana, * chwalcie imię Pańskie.\\
2. (ukłonić się) Niech będzie imię Pańskie błogosławione * odtąd i aż na wieki.\\
3. Od wschodu słońca aż do zachodu, * chwalebne jest imię Pańskie.\\
4. Wywyższony jest Pan nad wszystkie narody * i nad niebiosa chwała Jego.\\
5. Któż jest jak Pan, Bóg nasz, który mieszka na wysokości, * a na niskość spogląda na niebie i na ziemi?\\
6. Podnosi z ziemi biednego * i z gnoju wywyższa ubogiego:\\
7. Aby go posadzić z książętami, * z książętami ludu swego.\\
8. Który czyni, iż niepłodna mieszka w domu, * matka synów z weselem.\\
%---------------------
\red{\textbf{Ant.5}} Anioł Pański * ukazał się Józefowi mówiąc: Józefie, Synu Dawidów, nie bój się przyjąć Marji małżonki twej: albowiem co się w niej narodziło jest z Ducha Świętego: a porodzi syna i nazwiesz imię Jego Jezus.\\
\red{\textbf{Psalm 116}}\\
1. Chwalcie Pana wszystkie narody, * chwalcie Go wszystkie ludy:\\
2. Bo możne jest nad nami miłosierdzie Jego, * a prawda Pańska trwa na wieki.\\
%---------------------
\red{\textbf{Capitulum}}\\
Mąż wierny wielce chwalon będzie. A który strzeże Pana swego, uwielbiony. (Prz 28,20; 27,18)\\
%---------------------
\red{\textbf{Hymn}}\\
1. Ciebie, Józefie, sławią niebieskie zastępy, * Niechaj Twą chwałę głoszą także chrześcijan chóry * Co zasługami jaśniejący, czystym związkiem * Złączony jesteś z Najchwalebniejszą Dziewicą.\\
2. Zdumiony widzisz, że Twa Oblubienica ma zostać Matką. * Zwątpienie Cię przenika. * Anioł dopiero w widzeniu Ci objawia, * Że to Poczęcie tchnieniem Płomienia z wysoka.\\
3. Narodzonego Pana w objęcia przyjmujesz, * Uchodzisz z Nim w daleką krainę Egiptu, * W Jerozolimie zgubionego szukasz, płacząc, * I znajdujesz mieszając radość swą ze łzami.\\
4. Wszakże po zgonie wieczna szczęśliwość * Udziałem zwykłych śmiertelnych, śmierć palmę chwały im wręcza: * Ty już za życia, jak Święci, zażywasz Boga, * Przedziwnym losem jesteś tu od nich szczęśliwszy.\\
5. Najwyższa Trójco, przepuść kornie błagającym. * Daj dla zasług Józefa, do nieba się dostać. * Aby nam kiedyś, sługom Twym, wolno było, * Nieustający Tobie hymn wdzięczności śpiewać. * Amen.\\
%---------------------
\red{\Vbar.}~Postanowił go panem domu swego.\\
\red{\Rbar.}~I książęciem wszystkiej posiadłości swojej.\\
%---------------------
\red{\textbf{Ant.Magnificat}}~~A Józef, wstawszy * ze snu, uczynił jako mu rozkazał anioł; i przyjął żonę swoją.\\
\red{\textbf{Magnificat}}\\
1. Wielbij * duszo moja, Pana.\\
2. I rozradował się duch mój: * w Bogu, Zbawicielu moim.\\
3. Iż wejrzał na niskość służebnice swojéj: * albowiem oto odtąd błogosławioną mię zwać będą wszystkie narody.\\
4. Albowiem uczynił mi wielkie rzeczy, który możny jest: * i święte imię jego.\\
5. A miłosierdzie jego od narodu do narodów: * bojącym się jego.\\
6. Uczynił moc ramieniem swojem: * rozproszył pyszne myśli serca ich.\\
7. Złożył mocarze z stolice: * a podwyższył niskie.\\
8. Łaknące napełnił dobrami: * a bogacze z niczem puścił.\\
9. Przyjął Izraela, sługę swego: * wspomniawszy na miłosierdzie swoje.\\
10 Jako mówił do ojców naszych: * Abrahamowi i nasieniu jego na wieki.\\
%---------------------
\red{\textbf{Oratio}}\\
Módlmy się. Niechaj zasługi Oblubieńca Przenajświętszej Rodzicielki Twojej wspomogą, nas, prosimy Cię Panie, aby to, czego własnemi siłami otrzymać nie możemy, za jego wstawiennictwem było nam dane. Który żyjesz i królujesz z Bogiem Ojcem w jedności Ducha Świętego Bóg przez wszystkie wieki wieków. Amen.\\
\red{\textbf{Commemoratio}}\\
\red{\textbf{Ant.}}~~Albowiem jeden jest nauczyciel wasz, a wszyscy jesteście bracia. I ojca nie zówcie sobie na ziemi; albowiem jeden jest Ojciec wasz, który jest w niebiesiech. Ani się zówcie nauczycielmi, gdyż jeden nauczyciel wasz, Chrystus. (Mt 23,8-10)\\
\red{\Vbar.}~Aniołom swoim Bóg rozkazał o tobie. (Ps 90,11-12)\\
\red{\Rbar.}~Aby cię strzegli na wszystkich drogach twoich. \\ 
Módlmy się. Daj się przejednać, Panie, naszymi modłami i ulecz choroby dusz naszych, abyśmy otrzymawszy przebaczenie, zawsze się cieszyli Twoim błogosławieństwem. Przez naszego Pana Jezusa Chrystusa, Twojego Syna, który z Tobą żyje i króluje w jedności Ducha Świętego, Bóg, przez wszystkie wieki wieków.
\footnotetext{Tłumaczenia pochodzą ze strony divinumofficium.com oraz z Biblii w tł. ks. Jakuba Wujka SJ}